\chapter{总结与展望}

\section{研究总结}

本文围绕“基于大模型与软件运行日志的运维故障诊断方法”这一主题展开研究，核心关注的问题是：如何在大模型诊断框架下，围绕软件运行日志组织诊断输入，并在具体分析过程中以日志为主、结合监控指标辅助分析，从而提升软件运维故障诊断的稳定性与可追溯性。围绕这一问题，本文在梳理日志分析、监控指标与根因定位、多源观测协同、知识增强和结构化推理等相关研究的基础上，将研究对象收束为以日志分析为主、监控指标辅助的诊断输入，并在此基础上讨论知识增强与结构化推理的补充作用。

在方法层面，本文提出了CJW-Log方法。该方法并不将日志、指标、知识检索和结构化输出视为彼此并列的独立技术点，而是以日志分析为主链路，并结合监控指标辅助识别异常趋势、交叉印证日志线索，由大语言模型完成诊断推断；检索增强仅在知识依赖型场景中补充运维文档和案例知识，结构化推理则主要用于约束分析步骤和结果表达。与此对应，本文构建了统一的研究验证环境，在同一模型条件下完成受控对照验证与补充验证，从而保证不同设置之间的结果差异主要对应于输入组织方式与补充机制差异，而非执行框架差异。

验证结果表明，在本文受控条件下，以日志为主并结合监控指标辅助分析的输入组织构成诊断性能提升的主要来源。主实验中，Log+Metric 组的 Top-1 根因定位准确率达到 76.2\%，明显高于单源输入组；在此基础上加入知识增强后，Top-1 准确率进一步提高到 85.7\%，且增益主要集中在数据库连接池耗尽、慢 SQL 等知识依赖型场景；继续加入结构化推理后，Top-1 准确率未进一步提升，但输出规范率由 66.7\% 提高到 95.2\%，关键信息显式程度也得到改善。补充验证同样说明，结构化输出在跨数据集条件下对字段完整性和依据显式化具有稳定的正向作用。

需要说明的是，本文结论主要建立在受控验证环境之上，其作用在于说明CJW-Log方法在本文样例条件下的有效性。结合已有生产事故和微服务故障诊断研究可以看到，事故相关日志、指标、数据库信息以及Trace、拓扑等多源观测数据，都会影响根因定位的充分性~\cite{roy2024agents,zhang2024fdsurvey}。因此，本文关于日志与监控指标协同分析的结果，以及对跨服务传播型故障观测不足的讨论，能够与已有研究形成相互补充。真实生产环境中的高噪声日志、并发依赖、未知故障模式和多版本模型调用仍可能影响诊断稳定性，后续可在更大规模和更接近生产条件的数据上继续验证。

\section{未来展望}

后续研究可首先围绕复杂故障传播场景引入更丰富的观测信息。当前结果已经表明，跨服务传播型故障往往并不缺少“存在异常”的线索，而是缺少足以刻画传播路径的观测维度。因此，后续可从三个方向继续扩展：一是引入OpenTelemetry、Jaeger等分布式追踪数据，将日志、指标和Trace纳入统一分析上下文；二是在CJW-Log主链路中补充传播路径推断模块，把根因推断从单点定位扩展到路径级定位；三是结合CausalRCA、PC算法等因果发现方法，对服务依赖图和异常传播方向形成先验约束，再由大语言模型完成语义解释和排障建议生成~\cite{xin2023causalrca,ikram2022causaldiscovery,cheng2023rcasurvey}。

另一方面，仍有必要在更大规模样例和更接近真实环境的数据条件下继续开展验证，并同步完善知识库更新与筛选机制。随着样例规模、噪声复杂度和知识来源异构性增加，方法在稳定性、可解释性和检索鲁棒性上的表现仍有进一步评估空间。
